\documentclass[11pt,a4paper]{article}
\usepackage[utf8]{inputenc}
\usepackage[T1]{fontenc}
\usepackage[catalan]{babel}
\usepackage{lmodern}
\usepackage[sfdefault,lining]{FiraSans}
\usepackage[a4paper,top=2.3cm,bottom=2.2cm,left=1.8cm,right=1.8cm,headheight=1.3cm,headsep=0.4cm,footskip=1cm]{geometry}
\usepackage{amsmath,amssymb}
\usepackage{icomma}
\usepackage{graphicx}
\usepackage[table]{xcolor}
\usepackage{tikz}
\usepackage{fontawesome5}
\usepackage{ccicons}
\usepackage{fancyhdr}
\usepackage{lastpage}
\usepackage{tasks}
\usepackage{needspace}
\usepackage[most]{tcolorbox}
\usepackage{array}
\usepackage[hidelinks,pdftitle={Fitxa d'exercicis de la unitat 5: Geometria analítica en el pla. 1r de Batxillerat MI},pdfsubject={Matemàtiques I},pdfkeywords={geometria analítica, vectors, producte escalar, equacions de la recta, posicions relatives, rectes perpendiculars, angles, distàncies, mediatriu, bisectriu}]{hyperref}

\definecolor{navy}{HTML}{1B2748}
\definecolor{accent}{HTML}{E8603A}
\definecolor{teal}{HTML}{1F9E90}
\definecolor{soft}{HTML}{EEF1F8}
\definecolor{muted}{HTML}{5D6685}

\setlength{\parindent}{0pt}
\setlength{\parskip}{3pt}
\settasks{label=\alph*), label-format=\bfseries\color{navy}, label-width=1.3em, item-indent=1.8em, column-sep=0.6em, after-item-skip=5pt, before-skip=3pt, after-skip=4pt}

% ---------- capçalera i peu ----------
\newcommand{\llicencia}{\href{https://creativecommons.org/licenses/by-nc-sa/4.0/deed.ca}{\ccbyncsa}\ \ Llicència Creative Commons Reconeixement-NoComercial-CompartirIgual 4.0 Internacional (CC BY-NC-SA 4.0)}
\pagestyle{fancy}
\fancyhf{}
\fancyhead[L]{\raisebox{-0.35\height}{\includegraphics[height=0.95cm]{logo-gva.png}}}
\fancyhead[C]{\small\color{muted}Fitxa d'exercicis · Unitat 5: Geometria analítica en el pla · 1r Batxillerat MI}
\fancyhead[R]{\raisebox{-0.35\height}{\includegraphics[height=0.95cm]{logo-vermellar.png}}}
\fancyfoot[L]{\scriptsize\color{muted}\llicencia}
\fancyfoot[R]{\small\color{muted}\thepage}
\renewcommand{\headrulewidth}{0.4pt}
\renewcommand{\headrule}{\color{muted!40}\hrule height \headrulewidth}
\fancypagestyle{primera}{\fancyhf{}\renewcommand{\headrulewidth}{0pt}
  \fancyfoot[L]{\scriptsize\color{muted}\llicencia}
  \fancyfoot[R]{\small\color{muted}\thepage}}

% ---------- nivells i exercicis ----------
\newcommand{\estrelles}[1]{{\color{accent}\ifcase#1\or\faStar\faStar[regular]\faStar[regular]\or\faStar\faStar\faStar[regular]\or\faStar\faStar\faStar\fi}}
\newcounter{exer}
\newcommand{\ex}[1]{\par\needspace{4\baselineskip}\medskip\refstepcounter{exer}\noindent{\color{navy}\bfseries\theexer.}\ \ }

\newcommand{\seccio}[2]{\par\needspace{16\baselineskip}\bigskip
  \begin{tcolorbox}[colback=navy,colframe=navy,coltext=white,arc=2mm,boxrule=0pt,left=3mm,right=3mm,top=1.2mm,bottom=1.2mm]
  {\large\bfseries #1\quad #2}
  \end{tcolorbox}\vspace{-2pt}}

\newcommand{\exf}[1]{\par\medskip\refstepcounter{exer}\noindent{\color{navy}\bfseries\theexer.}\ \ }

% quadre de recordatori (només a la fitxa d'exercicis)
\newcommand{\recorda}[1]{\par\smallskip\begin{tcolorbox}[colback=accent!7,colframe=accent!55,arc=2mm,boxrule=0.5pt,left=2.5mm,right=2.5mm,top=1mm,bottom=1mm]\small{\color{accent}\faLightbulb[regular]}\ \textbf{Recorda:}\ #1\end{tcolorbox}}
\newcommand{\res}[1]{#1}\newcommand{\rest}[1]{#1}

\def\fmneed{11}% que el títol d'una secció no quede sol al final del full
% =====================================================================
% Capçalera FentMates de les fitxes: un color per a cada curs.
% Els colors són els de la pàgina inicial (LEVELS a assets/js/materials.js);
% el logo ve de _eines/logo/pdf/ (python3 _eines/logo/genera.py).
%
% Ús, al final del preàmbul de la unitat:
%   % =====================================================================
% Capçalera FentMates de les fitxes: un color per a cada curs.
% Els colors són els de la pàgina inicial (LEVELS a assets/js/materials.js);
% el logo ve de _eines/logo/pdf/ (python3 _eines/logo/genera.py).
%
% Ús, al final del preàmbul de la unitat:
%   % =====================================================================
% Capçalera FentMates de les fitxes: un color per a cada curs.
% Els colors són els de la pàgina inicial (LEVELS a assets/js/materials.js);
% el logo ve de _eines/logo/pdf/ (python3 _eines/logo/genera.py).
%
% Ús, al final del preàmbul de la unitat:
%   \input{../comu/fentmates.tex}
%   \fmcurs{1eso}
% i, al començament del document:
%   \capcalera{Matemàtiques}{1r ESO}{Fitxa d'exercicis · Unitat 2}{Divisibilitat}
% Opcional abans de \input: \def\fmneed{14} (espai mínim abans de cada secció).
% =====================================================================
\makeatletter
\def\fm@logodir{../../logo/pdf}
\providecommand{\fmneed}{10}
\def\fm@cursnom{}

% clau del curs, color i color del text damunt del color
\newcommand{\fm@curs}[3]{\@namedef{fm@color@#1}{#2}\@namedef{fm@text@#1}{#3}}
\fm@curs{1eso}{E6B84A}{navy}
\fm@curs{2eso}{3B82F6}{white}
\fm@curs{3eso}{FF7A50}{white}
\fm@curs{4eso-a}{A78BFA}{white}
\fm@curs{4eso-b}{6D28D9}{white}
\fm@curs{1batx-mi}{9B1C1C}{white}
\fm@curs{1batx-ccss}{EF6B6B}{white}
\fm@curs{2batx-mii}{1E6B47}{white}
\fm@curs{2batx-ccss}{2BB3A3}{white}

\newcommand{\fmcurs}[1]{%
  \@ifundefined{fm@color@#1}{\PackageError{fentmates}{Curs desconegut: #1}{}}{}%
  \edef\fm@tmp{\noexpand\definecolor{curs}{HTML}{\@nameuse{fm@color@#1}}%
               \noexpand\colorlet{curstext}{\@nameuse{fm@text@#1}}}\fm@tmp
  \colorlet{cursfosc}{curs!72!black}%
  \def\fm@logo{\fm@logodir/vertical-#1.pdf}%
  \def\fm@marca{\fm@logodir/marca-#1.pdf}}

% ---------- capçalera de la primera pàgina ----------
% #1 matèria (petit), #2 curs, #3 tipus de document i unitat, #4 títol
\newcommand{\capcalera}[4]{%
  \gdef\fm@cursnom{#2}%
  \par\noindent
  \begin{tcolorbox}[colback=curs,colframe=curs,boxrule=0pt,arc=2.2mm,width=3.5cm,height=2.15cm,
      valign=center,left=3mm,right=2mm,top=1mm,bottom=1mm,nobeforeafter,box align=center]
    \color{curstext}{\footnotesize #1}\par\vspace{1pt}{\fontsize{19}{21}\selectfont\bfseries #2\par}%
  \end{tcolorbox}\hspace{5mm}%
  \begin{minipage}[c]{\dimexpr\linewidth-3.5cm-5mm-2.9cm\relax}
    {\color{cursfosc}\small\bfseries\MakeUppercase{#3}}\par\vspace{2pt}
    {\color{navy}\fontsize{20}{23}\selectfont\bfseries #4\par}
  \end{minipage}\hfill
  \begin{minipage}[c]{2.6cm}\raggedleft\includegraphics[width=2.6cm]{\fm@logo}\end{minipage}\par
  \vspace{2pt}}

% ---------- barres de secció del color del curs ----------
\renewcommand{\seccio}[2]{\par\needspace{\fmneed\baselineskip}\bigskip
  \begin{tcolorbox}[colback=curs,colframe=curs,coltext=curstext,arc=2mm,boxrule=0pt,left=3mm,right=3mm,top=1.2mm,bottom=1.2mm]
  {\large\bfseries #1\quad #2}
  \end{tcolorbox}\vspace{-2pt}}

% ---------- capçalera i peu de les pàgines ----------
% primera pàgina: logos institucionals al peu
\fancypagestyle{primera}{\fancyhf{}\renewcommand{\headrulewidth}{0pt}
  \fancyfoot[L]{\raisebox{-0.32\height}{\includegraphics[height=0.78cm]{logo-gva.png}}\hspace{5mm}%
    \raisebox{-0.32\height}{\includegraphics[height=0.78cm]{logo-vermellar.png}}}
  \fancyfoot[R]{\scriptsize\color{muted}\href{https://creativecommons.org/licenses/by-nc-sa/4.0/deed.ca}{\ccbyncsa}\ \ CC BY-NC-SA 4.0 · fentmates.org\qquad\small\thepage}}
% la resta: el símbol a l'esquerra i el curs a la dreta, amb el títol al centre com abans
\fancyhead[L]{\raisebox{-0.3\height}{\includegraphics[height=0.62cm]{\fm@marca}}}
\fancyhead[R]{\small\bfseries\color{cursfosc}\fm@cursnom}
\renewcommand{\headrule}{\color{curs!45}\hrule height \headrulewidth}
\makeatother

%   \fmcurs{1eso}
% i, al començament del document:
%   \capcalera{Matemàtiques}{1r ESO}{Fitxa d'exercicis · Unitat 2}{Divisibilitat}
% Opcional abans de \input: \def\fmneed{14} (espai mínim abans de cada secció).
% =====================================================================
\makeatletter
\def\fm@logodir{../../logo/pdf}
\providecommand{\fmneed}{10}
\def\fm@cursnom{}

% clau del curs, color i color del text damunt del color
\newcommand{\fm@curs}[3]{\@namedef{fm@color@#1}{#2}\@namedef{fm@text@#1}{#3}}
\fm@curs{1eso}{E6B84A}{navy}
\fm@curs{2eso}{3B82F6}{white}
\fm@curs{3eso}{FF7A50}{white}
\fm@curs{4eso-a}{A78BFA}{white}
\fm@curs{4eso-b}{6D28D9}{white}
\fm@curs{1batx-mi}{9B1C1C}{white}
\fm@curs{1batx-ccss}{EF6B6B}{white}
\fm@curs{2batx-mii}{1E6B47}{white}
\fm@curs{2batx-ccss}{2BB3A3}{white}

\newcommand{\fmcurs}[1]{%
  \@ifundefined{fm@color@#1}{\PackageError{fentmates}{Curs desconegut: #1}{}}{}%
  \edef\fm@tmp{\noexpand\definecolor{curs}{HTML}{\@nameuse{fm@color@#1}}%
               \noexpand\colorlet{curstext}{\@nameuse{fm@text@#1}}}\fm@tmp
  \colorlet{cursfosc}{curs!72!black}%
  \def\fm@logo{\fm@logodir/vertical-#1.pdf}%
  \def\fm@marca{\fm@logodir/marca-#1.pdf}}

% ---------- capçalera de la primera pàgina ----------
% #1 matèria (petit), #2 curs, #3 tipus de document i unitat, #4 títol
\newcommand{\capcalera}[4]{%
  \gdef\fm@cursnom{#2}%
  \par\noindent
  \begin{tcolorbox}[colback=curs,colframe=curs,boxrule=0pt,arc=2.2mm,width=3.5cm,height=2.15cm,
      valign=center,left=3mm,right=2mm,top=1mm,bottom=1mm,nobeforeafter,box align=center]
    \color{curstext}{\footnotesize #1}\par\vspace{1pt}{\fontsize{19}{21}\selectfont\bfseries #2\par}%
  \end{tcolorbox}\hspace{5mm}%
  \begin{minipage}[c]{\dimexpr\linewidth-3.5cm-5mm-2.9cm\relax}
    {\color{cursfosc}\small\bfseries\MakeUppercase{#3}}\par\vspace{2pt}
    {\color{navy}\fontsize{20}{23}\selectfont\bfseries #4\par}
  \end{minipage}\hfill
  \begin{minipage}[c]{2.6cm}\raggedleft\includegraphics[width=2.6cm]{\fm@logo}\end{minipage}\par
  \vspace{2pt}}

% ---------- barres de secció del color del curs ----------
\renewcommand{\seccio}[2]{\par\needspace{\fmneed\baselineskip}\bigskip
  \begin{tcolorbox}[colback=curs,colframe=curs,coltext=curstext,arc=2mm,boxrule=0pt,left=3mm,right=3mm,top=1.2mm,bottom=1.2mm]
  {\large\bfseries #1\quad #2}
  \end{tcolorbox}\vspace{-2pt}}

% ---------- capçalera i peu de les pàgines ----------
% primera pàgina: logos institucionals al peu
\fancypagestyle{primera}{\fancyhf{}\renewcommand{\headrulewidth}{0pt}
  \fancyfoot[L]{\raisebox{-0.32\height}{\includegraphics[height=0.78cm]{logo-gva.png}}\hspace{5mm}%
    \raisebox{-0.32\height}{\includegraphics[height=0.78cm]{logo-vermellar.png}}}
  \fancyfoot[R]{\scriptsize\color{muted}\href{https://creativecommons.org/licenses/by-nc-sa/4.0/deed.ca}{\ccbyncsa}\ \ CC BY-NC-SA 4.0 · fentmates.org\qquad\small\thepage}}
% la resta: el símbol a l'esquerra i el curs a la dreta, amb el títol al centre com abans
\fancyhead[L]{\raisebox{-0.3\height}{\includegraphics[height=0.62cm]{\fm@marca}}}
\fancyhead[R]{\small\bfseries\color{cursfosc}\fm@cursnom}
\renewcommand{\headrule}{\color{curs!45}\hrule height \headrulewidth}
\makeatother

%   \fmcurs{1eso}
% i, al començament del document:
%   \capcalera{Matemàtiques}{1r ESO}{Fitxa d'exercicis · Unitat 2}{Divisibilitat}
% Opcional abans de \input: \def\fmneed{14} (espai mínim abans de cada secció).
% =====================================================================
\makeatletter
\def\fm@logodir{../../logo/pdf}
\providecommand{\fmneed}{10}
\def\fm@cursnom{}

% clau del curs, color i color del text damunt del color
\newcommand{\fm@curs}[3]{\@namedef{fm@color@#1}{#2}\@namedef{fm@text@#1}{#3}}
\fm@curs{1eso}{E6B84A}{navy}
\fm@curs{2eso}{3B82F6}{white}
\fm@curs{3eso}{FF7A50}{white}
\fm@curs{4eso-a}{A78BFA}{white}
\fm@curs{4eso-b}{6D28D9}{white}
\fm@curs{1batx-mi}{9B1C1C}{white}
\fm@curs{1batx-ccss}{EF6B6B}{white}
\fm@curs{2batx-mii}{1E6B47}{white}
\fm@curs{2batx-ccss}{2BB3A3}{white}

\newcommand{\fmcurs}[1]{%
  \@ifundefined{fm@color@#1}{\PackageError{fentmates}{Curs desconegut: #1}{}}{}%
  \edef\fm@tmp{\noexpand\definecolor{curs}{HTML}{\@nameuse{fm@color@#1}}%
               \noexpand\colorlet{curstext}{\@nameuse{fm@text@#1}}}\fm@tmp
  \colorlet{cursfosc}{curs!72!black}%
  \def\fm@logo{\fm@logodir/vertical-#1.pdf}%
  \def\fm@marca{\fm@logodir/marca-#1.pdf}}

% ---------- capçalera de la primera pàgina ----------
% #1 matèria (petit), #2 curs, #3 tipus de document i unitat, #4 títol
\newcommand{\capcalera}[4]{%
  \gdef\fm@cursnom{#2}%
  \par\noindent
  \begin{tcolorbox}[colback=curs,colframe=curs,boxrule=0pt,arc=2.2mm,width=3.5cm,height=2.15cm,
      valign=center,left=3mm,right=2mm,top=1mm,bottom=1mm,nobeforeafter,box align=center]
    \color{curstext}{\footnotesize #1}\par\vspace{1pt}{\fontsize{19}{21}\selectfont\bfseries #2\par}%
  \end{tcolorbox}\hspace{5mm}%
  \begin{minipage}[c]{\dimexpr\linewidth-3.5cm-5mm-2.9cm\relax}
    {\color{cursfosc}\small\bfseries\MakeUppercase{#3}}\par\vspace{2pt}
    {\color{navy}\fontsize{20}{23}\selectfont\bfseries #4\par}
  \end{minipage}\hfill
  \begin{minipage}[c]{2.6cm}\raggedleft\includegraphics[width=2.6cm]{\fm@logo}\end{minipage}\par
  \vspace{2pt}}

% ---------- barres de secció del color del curs ----------
\renewcommand{\seccio}[2]{\par\needspace{\fmneed\baselineskip}\bigskip
  \begin{tcolorbox}[colback=curs,colframe=curs,coltext=curstext,arc=2mm,boxrule=0pt,left=3mm,right=3mm,top=1.2mm,bottom=1.2mm]
  {\large\bfseries #1\quad #2}
  \end{tcolorbox}\vspace{-2pt}}

% ---------- capçalera i peu de les pàgines ----------
% primera pàgina: logos institucionals al peu
\fancypagestyle{primera}{\fancyhf{}\renewcommand{\headrulewidth}{0pt}
  \fancyfoot[L]{\raisebox{-0.32\height}{\includegraphics[height=0.78cm]{logo-gva.png}}\hspace{5mm}%
    \raisebox{-0.32\height}{\includegraphics[height=0.78cm]{logo-vermellar.png}}}
  \fancyfoot[R]{\scriptsize\color{muted}\href{https://creativecommons.org/licenses/by-nc-sa/4.0/deed.ca}{\ccbyncsa}\ \ CC BY-NC-SA 4.0 · fentmates.org\qquad\small\thepage}}
% la resta: el símbol a l'esquerra i el curs a la dreta, amb el títol al centre com abans
\fancyhead[L]{\raisebox{-0.3\height}{\includegraphics[height=0.62cm]{\fm@marca}}}
\fancyhead[R]{\small\bfseries\color{cursfosc}\fm@cursnom}
\renewcommand{\headrule}{\color{curs!45}\hrule height \headrulewidth}
\makeatother

\fmcurs{1batx-mi}

\begin{document}
\thispagestyle{primera}

\capcalera{Matemàtiques I}{1r Batx}{Fitxa d'exercicis · Unitat 5}{\raggedright Geometria analítica en el pla}
\vspace{10pt}
Nom i cognoms: \rule{8.4cm}{0.4pt}\hfill Grup: \rule{1.6cm}{0.4pt}\hfill Data: \rule{2.2cm}{0.4pt}
\vspace{8pt}

\begin{tcolorbox}[colback=soft,colframe=soft,arc=2mm,left=3mm,right=3mm,top=1.5mm,bottom=1.5mm]
\small Resol els exercicis al quadern mostrant tots els passos. Fes un dibuix aproximat en uns eixos de coordenades sempre que t'ajude: et permet detectar errors. Dona els resultats de manera exacta (fraccions i arrels) i, si cal, també aproximats. Les equacions de les rectes, en forma general $Ax+By+C=0$ amb coeficients enters, si no es diu el contrari.
\end{tcolorbox}

% =====================================================
\seccio{1}{Vectors i operacions amb vectors}

\exf{1} Donats els punts $A(3,\,-1)$ i $B(-2,\,4)$, calcula:\begin{tasks}(2)
  \task Les components de $\overrightarrow{AB}$ i de $\overrightarrow{BA}$.
  \task El mòdul de $\overrightarrow{AB}$.
  \task Un vector unitari amb la mateixa direcció i sentit que $\overrightarrow{AB}$.
  \task Un vector de mòdul 10 amb la mateixa direcció i sentit contrari.
\end{tasks}

\ex{1} Calcula l'argument (l'angle amb el semieix positiu $OX$) de cada vector. Fixa't en el quadrant!\begin{tasks}(3)
  \task $\vec a=(3,\,3)$
  \task $\vec a=(-1,\,1)$
  \task $\vec a=(-4,\,-4)$
  \task $\vec a=(0,\,-3)$
  \task $\vec a=(5,\,0)$
  \task $\vec a=(2,\,-2)$
\end{tasks}

\ex{1} Calcula les components d'un vector $\vec v$ si:\begin{tasks}(3)
  \task $|\vec v|=4$, argument $30^\circ$
  \task $|\vec v|=10$, argument $135^\circ$
  \task $|\vec v|=6$, argument $240^\circ$
\end{tasks}\recorda{$\vec v=\big(|\vec v|\cos\alpha,\ |\vec v|\sin\alpha\big)$.}

\ex{1} Siguen $\vec u=(2,\,-3)$, $\vec v=(-1,\,4)$ i $\vec w=(5,\,0)$. Calcula:\begin{tasks}(2)
  \task $\vec u+\vec v-\vec w$
  \task $3\vec u-2\vec v$
  \task $\tfrac12\vec w-\vec u+\vec v$
  \task $-2(\vec u-\vec v)+3\vec w$
\end{tasks}

\ex{2} Calcula $x$ i $y$ perquè es complisca $(x,\,3)+2(1,\,y)=(5,\,-1)$. Després, troba $k$ perquè el vector $(k,\,4)$ tinga mòdul 5.

\ex{2} Troba $m$ perquè els vectors $\vec u=(m,\,-2)$ i $\vec v=(6,\,3)$ siguen paral·lels. Per a aquest valor, quin nombre $k$ compleix $\vec u=k\vec v$?

\ex{2} Comprova que $\vec u=(1,\,2)$ i $\vec v=(3,\,-1)$ formen una base i escriu $\vec w=(1,\,9)$ com a combinació lineal de $\vec u$ i $\vec v$.


% =====================================================
\seccio{2}{Aplicacions geomètriques dels vectors}

\exf{1} Calcula els punts mitjans dels costats del triangle de vèrtexs $A(-2,\,1)$, $B(4,\,3)$ i $C(2,\,-5)$.

\ex{1} El punt $M(1,\,-2)$ és el punt mig del segment $AB$. Si $A(4,\,3)$, quines són les coordenades de $B$?

\ex{1} Troba els punts que divideixen el segment d'extrems $A(-3,\,4)$ i $B(9,\,-2)$ en tres parts iguals.

\ex{2} Divideix el segment d'extrems $A(1,\,-1)$ i $B(11,\,4)$ en cinc parts iguals.

\ex{1} Estudia si els punts estan alineats:\begin{tasks}(2)
  \task $A(1,\,3)$, $B(3,\,7)$, $C(-2,\,-3)$
  \task $A(0,\,2)$, $B(3,\,4)$, $C(6,\,7)$
\end{tasks}Calcula $k$ perquè $A(-1,\,2)$, $B(2,\,k)$ i $C(5,\,-4)$ estiguen alineats.

\ex{1} Calcula el simètric del punt $P(2,\,-3)$ respecte del punt $Q(-1,\,4)$, i el simètric de $R(5,\,1)$ respecte de l'origen.

\ex{2} Tres vèrtexs consecutius d'un paral·lelogram $ABCD$ són $A(-1,\,-2)$, $B(4,\,-1)$ i $C(5,\,3)$. Troba el vèrtex $D$ i el centre del paral·lelogram.


% =====================================================
\seccio{3}{Producte escalar i angle de dos vectors}

\exf{1} Siguen $\vec u=(3,\,-2)$ i $\vec v=(4,\,1)$. Calcula $\vec u\cdot\vec v$, $|\vec u|$, $|\vec v|$ i l'angle que formen.

\ex{1} Calcula l'angle que formen els vectors:\begin{tasks}(2)
  \task $(1,\,2)$ i $(-2,\,1)$
  \task $(1,\,1)$ i $(-1,\,0)$
  \task $(3,\,4)$ i $(-4,\,-3)$
  \task $(2,\,-1)$ i $(1,\,3)$
\end{tasks}

\ex{2} Si $|\vec u|=2$, $|\vec v|=5$ i l'angle que formen és de $120^\circ$, calcula:\begin{tasks}(3)
  \task $\vec u\cdot\vec v$
  \task $(2\vec u+\vec v)\cdot(\vec u-\vec v)$
  \task $|\vec u+\vec v|$
\end{tasks}\recorda{$\vec u\cdot\vec u=|\vec u|^2$ i el producte escalar és distributiu.}

\ex{2} Troba $m$ perquè $\vec u=(m,\,2)$ i $\vec v=(3,\,-6)$ siguen perpendiculars. Troba també $k$ perquè $(1,\,k)$ forme un angle de $60^\circ$ amb $(1,\,0)$.

\ex{1} Troba els vectors unitaris perpendiculars a $\vec v=(6,\,-8)$.

\ex{2} Comprova que el triangle de vèrtexs $A(1,\,1)$, $B(5,\,3)$ i $C(3,\,7)$ és rectangle i isòsceles. Quant mesuren els seus angles?

\ex{2} Calcula la projecció del vector $\vec u=(2,\,5)$ sobre la direcció de $\vec v=(3,\,-4)$. Què indica el signe?


% =====================================================
\seccio{4}{Equacions de la recta}

\exf{1} Escriu totes les equacions (vectorial, paramètriques, contínua, punt-pendent, explícita i general) de la recta que passa per $P(2,\,-1)$ i té vector director $\vec u=(3,\,4)$.

\ex{1} Escriu totes les equacions de la recta que passa per $A(-1,\,3)$ i $B(3,\,-5)$.

\ex{1} Troba les equacions punt-pendent, explícita i general de la recta que passa per $P(4,\,-2)$ i té pendent $-\tfrac32$.

\ex{1} Donada la recta $r:2x-5y+10=0$, troba: el pendent, els punts de tall amb els eixos, un vector director, un vector normal i unes equacions paramètriques.

\ex{1} Comprova si els punts $P(7,\,1)$ i $Q(4,\,0)$ pertanyen a la recta $r:\left\{\begin{array}{l}x=1+3\lambda\\ y=-3+2\lambda\end{array}\right.$ Calcula $k$ perquè $(k,\,5)$ siga de $r$.

\ex{1} Escriu l'equació de la recta:\begin{tasks}(2)
  \task Paral·lela a l'eix $OX$ per $(3,\,-2)$.
  \task Paral·lela a l'eix $OY$ per $(3,\,-2)$.
  \task Que passa per l'origen i per $(3,\,-2)$.
  \task Bisectriu del 1r i 3r quadrants.
\end{tasks}

\ex{1} Troba l'equació general de la recta que passa per $P(1,\,5)$ i té vector normal $\vec n=(3,\,-1)$.

\ex{2} Troba l'equació de la recta que passa per $(2,\,3)$ i forma un angle de $135^\circ$ amb el semieix positiu $OX$, i la de la recta que talla l'eix $OY$ en $(0,\,-4)$ i forma un angle de $60^\circ$ amb $OX$.


% =====================================================
\seccio{5}{Posicions relatives de dues rectes}

\exf{1} Estudia la posició relativa de cada parell de rectes i, si es tallen, troba el punt de tall:\begin{tasks}(1)
  \task $r:2x-3y+1=0$, \quad $s:4x+y-5=0$
  \task $r:x-2y+3=0$, \quad $s:-3x+6y-9=0$
  \task $r:3x+y-2=0$, \quad $s:6x+2y+1=0$
\end{tasks}

\ex{1} Estudia la posició relativa de $r:\left\{\begin{array}{l}x=1+2\lambda\\ y=-1+\lambda\end{array}\right.$ i $s:\left\{\begin{array}{l}x=3-4\mu\\ y=2\mu\end{array}\right.$

\ex{1} Estudia la posició relativa de $r:\left\{\begin{array}{l}x=2+\lambda\\ y=1-3\lambda\end{array}\right.$ i $s:3x+y-7=0$.

\ex{2} Troba $a$ perquè les rectes $ax+3y-2=0$ i $2x-y+5=0$ siguen: a) secants; b) paral·leles. Poden ser coincidents?

\ex{2} Calcula $m$ i $n$ perquè les rectes $mx+2y+n=0$ i $3x-y+4=0$ siguen coincidents.

\ex{2} Les rectes $r:x-y+1=0$, $s:x+y-5=0$ i $t:y=-1$ determinen un triangle. Troba'n els vèrtexs.


% =====================================================
\seccio{6}{Paral·lelisme i perpendicularitat}

\exf{1} Troba les rectes paral·lela i perpendicular a $r:3x-2y+5=0$ que passen per $P(2,\,-1)$.

\ex{1} Troba les rectes paral·lela i perpendicular a $r:\left\{\begin{array}{l}x=1-2\lambda\\ y=3+\lambda\end{array}\right.$ que passen per $A(0,\,-2)$. Dona-les en forma general.

\ex{1} Donada $r:y=-\tfrac13x+2$, troba les equacions explícites de les rectes paral·lela i perpendicular a $r$ que passen per $(-3,\,4)$.

\ex{2} Calcula $k$ perquè les rectes $kx-2y+1=0$ i $3x+(k+1)y-4=0$ siguen perpendiculars. Poden ser paral·leles?

\ex{2} En el triangle de vèrtexs $A(1,\,1)$, $B(7,\,3)$ i $C(3,\,7)$, troba l'equació de l'altura que surt de $C$ i el peu d'aquesta altura.

\ex{2} Troba la mediatriu del segment d'extrems $A(-3,\,2)$ i $B(5,\,-4)$. Comprova-la amb la definició (punts que equidisten de $A$ i de $B$).


% =====================================================
\seccio{7}{Angle de dues rectes}

\exf{1} Calcula l'angle que formen les rectes $r:3x-y+2=0$ i $s:x+2y-1=0$.

\ex{1} Calcula l'angle de $r:\left\{\begin{array}{l}x=1+2\lambda\\ y=3-\lambda\end{array}\right.$ i $s:\left\{\begin{array}{l}x=2-\lambda\\ y=4+3\lambda\end{array}\right.$

\ex{1} Calcula l'angle de les rectes $y=2x-1$ i $y=-x+4$.

\ex{1} Quin angle forma cada recta amb l'eix $OX$? \quad a) $3x-3y+1=0$ \quad b) $x+2y=0$

\ex{3} Troba les rectes que passen per $P(1,\,2)$ i formen un angle de $45^\circ$ amb la recta $r:x-3y+1=0$.


% =====================================================
\seccio{8}{Distàncies}

\exf{1} Calcula les longituds dels costats del triangle de vèrtexs $A(-1,\,2)$, $B(3,\,5)$ i $C(6,\,1)$ i classifica'l. Quin és el seu perímetre?

\ex{1} Calcula la distància del punt $P(3,\,-2)$ a cadascuna de les rectes:\begin{tasks}(2)
  \task $4x-3y+2=0$
  \task $y=2x+1$
  \task $x=-1$
  \task $\left\{\begin{array}{l}x=1+\lambda\\ y=2-\lambda\end{array}\right.$
\end{tasks}

\ex{1} Comprova que les rectes $r:3x+4y-5=0$ i $s:6x+8y+15=0$ són paral·leles i calcula la distància entre elles.

\ex{2} Troba els punts de l'eix $OY$ que estan a distància 2 de la recta $3x-4y+1=0$.

\ex{2} Troba les rectes paral·leles a $5x-12y+3=0$ que estan a distància 2 d'ella.

\ex{2} Una circumferència té el centre en $C(2,\,-1)$ i és tangent a la recta $3x+4y-12=0$. Quin és el seu radi?


% =====================================================
\seccio{9}{Problemes mètrics}

\exf{2} Troba la projecció del punt $P(-3,\,4)$ sobre la recta $r:2x-y+5=0$ i el simètric de $P$ respecte de $r$.

\ex{1} Calcula el simètric del punt $P(4,\,1)$ respecte de la recta $y=x$ i respecte de la recta $x=2$.

\ex{3} Troba les bisectrius dels angles que formen les rectes $r:3x+4y-1=0$ i $s:5x-12y+2=0$. Comprova que són perpendiculars.

\ex{2} Calcula l'àrea del triangle determinat per les rectes $r:x-y+1=0$, $s:x+y-5=0$ i $t:y=-1$ (exercici 35).

\ex{2} En el triangle $A(-2,\,0)$, $B(4,\,2)$, $C(1,\,6)$, calcula:\begin{tasks}(2)
  \task La mitjana que surt de $A$.
  \task L'altura que surt de $A$.
  \task El baricentre.
  \task L'àrea.
\end{tasks}

\ex{2} Troba el circumcentre i el radi de la circumferència circumscrita al triangle $A(0,\,0)$, $B(6,\,0)$, $C(0,\,4)$.

\ex{2} Troba l'ortocentre del triangle de vèrtexs $A(-3,\,0)$, $B(3,\,0)$ i $C(1,\,4)$.

\ex{2} Els punts $A(1,\,1)$ i $C(5,\,5)$ són vèrtexs oposats d'un quadrat. Troba els altres dos vèrtexs i l'àrea del quadrat.

\ex{3} Troba el punt de la recta $y=x$ que està a la mateixa distància de $A(1,\,0)$ que de $B(5,\,2)$.

\vfill
\begin{center}\color{muted}\small\faCheckCircle[regular]\ \ Ànim! Comprova els teus resultats amb les eines interactives dels apunts de la unitat.\end{center}

\end{document}
